\documentclass{article}

\usepackage[T1]{fontenc}
\usepackage{array}
\usepackage{geometry}
\usepackage{hyperref}

\setlength{\parindent}{0pt}
\setlength{\parskip}{0.6em}

\geometry{margin=1in}

\title{Simulated Market Agents and Instrument Behaviour}
\author{Will Lucas}
\date{}

\begin{document}

\maketitle
\tableofcontents

\section{Purpose}

This document describes the intended high-level behaviour of automated
trading agents and simulated instruments.

The aim is to create an active two-sided market in which traded
prices emerge from orders submitted by participants. The exchange itself does
not arbitrarily change a displayed market price. It accepts orders, applies
price--time priority, and records the prices at which trades occur.

\section{Instrument Behaviour}

\subsection{Fundamental value}

Each instrument has an underlying fundamental value representing an estimate
of its economic worth. This value exists separately from the last traded price
and from the prices of orders currently resting on the book.

The fundamental value changes periodically as simulation time advances. It is
a driver of agent behaviour rather than an execution price imposed by the
exchange. Consequently, the traded price may be above or below fundamental
value, and it may take time for market activity to move the two closer
together.

When an instrument is first issued, its initial fundamental value is the
configured starting price. The startup configuration distributes issued volume
across smaller sell orders around that price. The instrument begins in a neutral sentiment
regime. Its initial volatility is 0.1 per cent. These defaults provide a
complete hidden state without requiring additional instrument-creation
parameters.

\subsection{Sentiment}

Each instrument has a sentiment regime, such as bullish, neutral,
or bearish. Sentiment influences the likely direction of changes in its
fundamental value:

\begin{itemize}
    \item bullish sentiment tends to produce positive changes;
    \item neutral sentiment tends to produce little directional movement; and
    \item bearish sentiment tends to produce negative changes.
\end{itemize}

Sentiment should normally persist for a period before changing. This allows
recognisable trends to develop instead of selecting an unrelated direction on
every simulation step. Regime changes remain uncertain and may reverse an
existing trend.

On each simulation step there is a small possibility that an instrument moves
to a different sentiment regime. Most steps retain the current regime, allowing
market conditions to remain persistent.

Sentiment does not mean that every agent automatically buys or sells in the
same direction. It affects the underlying environment, while agents interpret
available information according to their own strategies.

\subsection{Volatility and random movement}

Volatility controls the typical size and variability of changes in fundamental
value. A low-volatility instrument should usually move gradually, while a
high-volatility instrument may experience larger and less predictable changes.

Random movement prevents the market from following a perfectly predictable
path. Sentiment supplies directional pressure, but volatility and random shocks
mean that a bullish instrument may fall on a particular step and a bearish
instrument may temporarily rise.

Volatility is represented as a proportion of fundamental value rather than a
fixed number of ticks. It fluctuates gradually on every simulation step within
bounded limits. Fundamental value then receives a random percentage movement,
with bullish or bearish sentiment biasing the likely direction. Fundamental
value is always kept at or above one tick.

The simulator does not require reproducible random outcomes. Separate runs may
produce completely different price histories even when they begin with the
same instruments and participants.

\section{Agent Categories}

\subsection{Noise traders}

Noise traders create ordinary, imperfect market activity. They submit buys and
sells with variation in side, price, quantity, and timing. Their decisions may
be loosely centred around public market prices, but they do not possess special
knowledge of the underlying fundamental value.

Noise traders help populate both sides of the order book and prevent all order
flow from being the product of highly informed strategies. Their activity may
temporarily move traded prices away from fundamental value.

\subsection{Fundamental traders}

Fundamental traders attempt to identify when the market price differs from an
instrument's economic worth. They may buy when the instrument appears cheap
and sell when it appears expensive.

They do not receive the exact hidden fundamental value or sentiment regime.
Instead, they receive an imperfect estimate or signal. Different traders may
receive signals of different quality or interpret the same signal differently.

\subsection{Trend followers}

Trend followers make decisions from observable market behaviour. They may
respond to recent trade prices, changes in price, momentum, trading volume, or
the state of the public order book.

They do not need access to fundamental information. A trend follower may buy
into an upward movement or sell into a downward movement even where the market
has already moved away from fundamental value. This can extend trends and
introduce delayed reactions to changing conditions.

\subsection{Market makers}

Market makers attempt to provide continuous liquidity by maintaining buy and
sell orders around an estimated fair price. They earn a potential spread in
exchange for accepting inventory risk and the risk of trading against better
informed participants.

Their decisions may consider the public order book, recent trades, current
inventory, available cash, and an imperfect estimate of fundamental value.
They should adjust their quotes or reduce activity when inventory becomes too
one-sided.

\section{Information Access}

Hidden market state should not be exposed equally to all agents. Public market
information is available to strategies that choose to use it, while
fundamental information is supplied only as an imperfect signal. No ordinary
agent receives the exact fundamental value or the literal hidden sentiment
regime.

\begin{center}
\small
\renewcommand{\arraystretch}{1.25}
\begin{tabular}{|>{\raggedright\arraybackslash}p{3.6cm}|c|c|c|c|}
\hline
\textbf{Information} & \textbf{Noise trader} & \textbf{Fundamental trader}
& \textbf{Trend follower} & \textbf{Market maker} \\
\hline
Public order book & Optional & Yes & Yes & Yes \\
\hline
Public trade history & Optional & Yes & Yes & Yes \\
\hline
Exact fundamental value & No & No & No & No \\
\hline
Imperfect fundamental signal & No & Yes & No & Optional \\
\hline
Hidden sentiment regime & No & Indirectly & No & Indirectly \\
\hline
Own cash, positions, and reservations & Yes & Yes & Yes & Yes \\
\hline
Other participants' private state & No & No & No & No \\
\hline
\end{tabular}
\end{center}

``Optional'' means that a strategy may use the information but does not depend
on it. ``Indirectly'' means that sentiment may influence an agent's imperfect
signal without revealing the name or exact state of the regime.

\section{Market Price Formation}

Agents convert their observations and strategies into order requests. Those
orders interact through the normal exchange matching rules. The resulting
trade prices form the public price history of each instrument.

This separation is important:

\begin{enumerate}
    \item the simulated environment changes an instrument's hidden fundamental
          conditions;
    \item agents receive different observations of the environment and market;
    \item agents submit orders based on those observations;
    \item the exchange matches compatible orders; and
    \item completed trades determine the observable market price.
\end{enumerate}

\end{document}
